\section{Half-return maps, inward graphs, tail balance, and hyperbolicity}\label{sec:toolkit}


The reflected half-returns satisfy the variational identities used
for trinomials by Chen, Jin, Li and Zhou. Their running-minimum
estimate controls contraction along an inward leg
\cite[Proposition~2.2, (3.5)--(3.7) and (3.17)--(3.19)]{CJLZ}.
We shall need these estimates both for finite turns and for inward
graphs with a prescribed limiting ordinate at infinity.


\subsection{Reflection, Filippov potential coordinates, and variational identities}\label{subsec:half-returns}
Given $F$ as in \eqref{eq:collected}, define the two half-profiles $P_+(r)=F(r)$ and $P_-(r)=F(-r)$
for $r\ge 0$. On the right half-plane $\{x\ge 0\}$ with $r=x$ and forward time $t$, and on the left
half-plane $\{x\le 0\}$ under the reflection $r=-x\ge 0$ and time reversal $\tilde{t}=-t$, system
\eqref{eq:lienard} takes the common form
\begin{equation}\label{eq:half-system}
  \dot{r} = y - P_\sigma(r), \qquad \dot{y} = -r \qquad (r\ge 0,\ \sigma\in\{+,-\}).
\end{equation}

Throughout $\{r>0\}$ one has $\dot{y}=-r<0$, so the ordinate $y$ decreases strictly along every
trajectory segment in $\{r>0\}$. Filippov introduced the potential separately on the two half-axes
\cite[pp.~171--172, (2)--(6)]{Filippov}. 
The potential coordinate $u=r^2/2\ge 0$
expresses each half-trajectory in $\{r>0\}$ as a single-valued smooth graph $u=u_\sigma(y)$ satisfying
\begin{equation}\label{eq:filippov}
  \frac{\dd u_\sigma}{\dd y} = \frac{r\dot{r}}{\dot{y}} = P_\sigma(\sqrt{2u_\sigma}) - y.
\end{equation}

At any interior critical point of $u_\sigma(y)$, where $y=P_\sigma(\sqrt{2u_\sigma})$, differentiating
\eqref{eq:filippov} with respect to $y$ yields $u_{\sigma,yy}=-1<0$. Consequently, each half-arc in
$\{r>0\}$ possesses at most one stationary radius, which is necessarily a strict global maximum
$r=t_\sigma>0$ called its \emph{turn}.


For $\beta>0$, launch the trajectory of \eqref{eq:half-system} from $(0,\beta)$ into $\{r>0\}$.
Whenever it first returns transversely to $\{r=0\}$ at a finite negative ordinate, denote that
ordinate by $b_\sigma(\beta)<0$, and define the half-return displacement
\[
  \oscD(\beta) = b_+(\beta) - b_-(\beta)
\]
on the common open domain of transverse returns. To analyze multiscale limits uniformly, given radial
and ordinate scales $\Lambda,M>0$ we set $\eps=(\Lambda/M)^2$, $w=y/M$, $\tilde{u}=u/\Lambda^2=(r/\Lambda)^2/2$, and
$\phi_\sigma(\tilde{u})=M^{-1}P_\sigma(\Lambda\sqrt{2\tilde{u}})$. For a half-arc turning at $\tilde{u}=T>0$
(physical radius $\Lambda\sqrt{2T}$) that meets both vertical half-axes transversely at $(0,Mw_o(0))$
($w_o(0)>0$) and $(0,Mw_i(0))$ ($w_i(0)<0$), its outward leg $w_o(\tilde{u})$ and inward leg
$w_i(\tilde{u})$ on $[0,T]$ solve
\begin{equation}\label{eq:scaled-legs}
  \frac{\dd w}{\dd\tilde{u}} = \frac{\eps}{\phi_\sigma(\tilde{u})-w}, \qquad w_o(\tilde{u}) >
  \phi_\sigma(\tilde{u}) > w_i(\tilde{u}) \quad (0<\tilde{u}<T), \qquad w_o(T)=w_i(T)=\phi_\sigma(T).
\end{equation}
Define the half-squared scaled axis intercepts $X_\sigma(T)=\frac{1}{2}w_o(0)^2$ and
$Y_\sigma(T)=\frac{1}{2}w_i(0)^2$, together with the suffix running extrema on $[0,T]$:
\begin{equation}\label{eq:running-extrema}
  m_{\sigma,T}(\tilde{u}) = \max_{s\in[\tilde{u},T]} \phi_\sigma(s), \qquad l_{\sigma,T}(\tilde{u}) =
  \min_{s\in[\tilde{u},T]} \phi_\sigma(s).
\end{equation}

\begin{lemma}[Single crossing, variational derivatives, and running-extremum bounds]\label{lem:variational}
Let $F$ be a real polynomial with $F(0)=0$.
\begin{enumerate}
  \item[\textup{(i)}] Every periodic orbit of \eqref{eq:lienard} surrounds $(0,0)$ and intersects each
  of the rays $\{x=0,y>0\}$ and $\{x=0,y<0\}$ transversely in a single point. Periodic orbits of
  \eqref{eq:lienard} stand in bijection with zeros $\beta>0$ of $\oscD(\beta)=b_+(\beta)-b_-(\beta)$,
  and at any matched orbit $\oscD(\beta_*)=0$ the first-return map $\Pi(\beta)=b_-^{-1}(b_+(\beta))$
  has Floquet multiplier $\Pi'(\beta_*)=b_+'(\beta_*)/b_-'(\beta_*)$.
  \item[\textup{(ii)}] On the domain of transverse returns, $b_\sigma(\beta)<0$ is real analytic and strictly decreasing in $\beta$, with
  \begin{equation}\label{eq:b-deriv}
    b_\sigma'(\beta) = \frac{\beta}{b_\sigma(\beta)} e^{-J_\sigma(\beta)}, \qquad J_\sigma(\beta) =
    \int_{\mathrm{arc}} P_\sigma'(r(t))\dd t = \int_{b_\sigma(\beta)}^{\beta}
    \frac{P_\sigma'(\sqrt{2u_\sigma(y)})}{\sqrt{2u_\sigma(y)}}\dd y.
  \end{equation}
  At a matched periodic orbit $\oscD(\beta_*)=0$, one has
  $\Pi'(\beta_*)=\exp(J_-(\beta_*)-J_+(\beta_*))$ and $\sgn(\Pi'(\beta_*)-1)=-\sgn\oscD'(\beta_*)$.
  Furthermore, if $P_{\sigma,\mu}(r)=P_{\sigma,0}(r)+\mu\,\partial_\mu P_{\sigma,0}(r)+o(\mu)$ varies
  smoothly with a parameter $\mu$ at fixed $\beta$, then
  \begin{equation}\label{eq:b-param}
    \partial_\mu b_\sigma(\beta,\mu) = -\frac{1}{b_\sigma(\beta,\mu)}
    \int_{b_\sigma(\beta,\mu)}^{\beta} \exp\!\left(-\int_{b_\sigma(\beta,\mu)}^{y}
    \frac{P_{\sigma,\mu}'(\sqrt{2u_\sigma(v)})}{\sqrt{2u_\sigma(v)}}\dd v\right) \partial_\mu
    P_{\sigma,\mu}(\sqrt{2u_\sigma(y)})\dd y.
  \end{equation}
  For two polynomial half-profiles $P_1,P_2$ vanishing at zero, fix a common
  launch $\beta>0$ for which both half-arcs return transversely to
  $\{r=0,y<0\}$. If $P_2\ge P_1$ on the union of their radial ranges,
  with strict inequality on a nonempty open radial interval traversed
  by both arcs, then $b_{P_2}(\beta)>b_{P_1}(\beta)$.
  \item[\textup{(iii)}] For any turn $T>0$ with transverse axis hits in the scaled coordinates
  \eqref{eq:scaled-legs}, $X_\sigma(T)$ and $Y_\sigma(T)$ increase strictly with $T$ and satisfy
  \begin{equation}\label{eq:turn-derivs}
    X_\sigma'(T) = \eps e^{J_{o,\sigma}(T)}, \quad Y_\sigma'(T) = \eps e^{-J_{i,\sigma}(T)}, \quad
    -\eps J_{o,\sigma}(T) = \int_0^T \phi_\sigma'\dd w_o, \quad \eps J_{i,\sigma}(T) = \int_0^T
    \phi_\sigma'\dd w_i,
  \end{equation}
  where $J_\sigma = J_{o,\sigma}+J_{i,\sigma}$. Moreover, for every $\tilde{u}\in[0,T]$,
  \begin{equation}\label{eq:envelope-bounds}
    0 \le w_o(\tilde{u}) - m_{\sigma,T}(\tilde{u}) \le \sqrt{2\eps(T-\tilde{u})}, \qquad 0 \le
    l_{\sigma,T}(\tilde{u}) - w_i(\tilde{u}) \le \sqrt{2\eps(T-\tilde{u})}.
  \end{equation}
\end{enumerate}
\end{lemma}


\begin{proof}
(i) Because $\dot{y}=-x$ vanishes only on $\{x=0\}$ and $(0,0)$ is the sole equilibrium of
\eqref{eq:lienard} (as $y-F(0)=y=0$ on $\{x=0\}$), no periodic orbit can lie in $\{x\ge 0\}$ or
$\{x\le 0\}$, and by index theory every periodic orbit $\Gamma$ is a simple closed Jordan curve
enclosing $(0,0)$. Along $\{x=0,y>0\}$ one has $\dot{x}=y>0$, so $\Gamma$ crosses $\{x=0,y>0\}$
transversely from left to right at every intersection. The signed intersection number with this ray is
the winding number of
the Jordan curve about the origin, of absolute value one. All crossings
have the same sign, so there is exactly one. On the negative ray all
crossings again have the same sign. Reversing time on $\{x\le 0\}$ gives $\oscD(\beta)=0$ as the exact
closure condition and $\Pi=b_-^{-1}\circ b_+$.

(ii) Near a transverse axis hit $y_0\ne 0$ ($\beta>0$ or $b_\sigma<0$), the planar trajectory
satisfies $y(r)=y_0-r^2/(2y_0)+O(r^3)$, hence $u_\sigma(y)=-y_0(y-y_0)+O(|y-y_0|^{3/2})$ and
$u_{\sigma,y}(y_0)=-y_0\ne 0$. Since $F(0)=0$, the coefficient
$g_\sigma(y)=P_\sigma'(\sqrt{2u_\sigma(y)})/\sqrt{2u_\sigma(y)}$ is
$O(u_\sigma(y)^{-1/2})=O(|y-y_0|^{-1/2})$, which is integrable at both endpoints $b_\sigma$ and
$\beta$. Differentiating \eqref{eq:filippov} and $u_\sigma(\beta,\beta,\mu)=0$ with respect to
$(\beta,\mu)$ gives the linear variational equation $\partial_y(\delta u_\sigma)=g_\sigma(y)\,\delta
u_\sigma + \partial_\mu P_{\sigma,\mu}(\sqrt{2u_\sigma})\,\delta\mu$ with boundary value $\delta
u_\sigma(\beta)=\beta\,\delta\beta$ (since $u_{\sigma,y}(\beta)\,\delta\beta+\delta u_\sigma(\beta)=0$
and $u_{\sigma,y}(\beta)=-\beta$). Integrating from $\beta$ down to $b_\sigma$ yields
\[
  \delta u_\sigma(b_\sigma) = \beta e^{-J_\sigma(\beta)}\,\delta\beta - \left(\int_{b_\sigma}^\beta
  e^{-\int_{b_\sigma}^y g_\sigma(v)\dd v}\,\partial_\mu P_{\sigma,\mu}(\sqrt{2u_\sigma(y)})\dd
  y\right)\delta\mu,
\]
where $\int_{b_\sigma}^\beta g_\sigma(y)\dd y=\int_{\mathrm{arc}}P_\sigma'(r(t))\dd t=J_\sigma(\beta)$
via $\dd y=-r\dd t=-\sqrt{2u_\sigma}\dd t$. Differentiating
$u_\sigma(b_\sigma(\beta,\mu),\beta,\mu)=0$ gives $-b_\sigma\,\delta b_\sigma+\delta
u_\sigma(b_\sigma)=0$, proving \eqref{eq:b-deriv} and \eqref{eq:b-param}. Because $b_-'(\beta_*)<0$,
$\Pi'(\beta_*)-1=\oscD'(\beta_*)/b_-'(\beta_*)$ has sign $-\sgn\oscD'(\beta_*)$. For the finite comparison write
$b_i=b_{P_i}(\beta)$, $r_i=\sqrt{2u_i}$ and compare the two graphs
at common ordinates $\max(b_1,b_2)\le y\le\beta$. Their difference
$d=u_2-u_1$ solves
\[
 d'=g(y)d+P_2(r_1)-P_1(r_1),\qquad d(\beta)=0,
\]
where $g$ is the divided difference of $u\mapsto P_2(\sqrt{2u})$.
In the interior $|g|\le 2\sup|P_2'|/(r_1+r_2)$, with the supremum
over the two radial ranges. The transverse expansions above make
this bound integrable even at a shared hit. Integrating backwards
gives $d\le0$. If $b_2\le b_1$, the common ordinate interval contains
the whole first excursion. There the forcing is positive on an open
set, hence $d(b_1)<0$, whereas $d(b_1)=u_2(b_1)\ge0$. Thus $b_2>b_1$.
Two returning arcs of the same
half-field enclose a compact strip with the axis. Intermediate launches
return in that strip by uniqueness and monotonicity of $y$.

(iii) In scaled variables \eqref{eq:scaled-legs}, view $\tilde{u}=\tilde{u}_o(w,T)$ on
$[\phi_\sigma(T),w_o(0)]$ as solving $\eps\,\partial_w\tilde{u}_o=\phi_\sigma(\tilde{u}_o)-w$ with
$\tilde{u}_o(\phi_\sigma(T),T)=T$. Since $\partial_w\tilde{u}_o(\phi_\sigma(T),T)=0$, differentiating
at $w=\phi_\sigma(T)$ gives $(\partial_T\tilde{u}_o)(\phi_\sigma(T),T)=1$. Integrating
$\eps\,\partial_w(\partial_T\tilde{u}_o)=\phi_\sigma'(\tilde{u}_o)\,\partial_T\tilde{u}_o$ from
$\phi_\sigma(T)$ to $w_o(0)$ gives
\[
  (\partial_T\tilde{u}_o)\bigl(w_o(0),T\bigr) =
  \exp\!\left(\frac{1}{\eps}\int_{\phi_\sigma(T)}^{w_o(0)}\phi_\sigma'\bigl(\tilde{u}_o(w)\bigr)\dd
  w\right) = e^{J_{o,\sigma}(T)},
\]
where $-\eps J_{o,\sigma}(T)=\int_0^T\phi_\sigma'(\tilde{u})\dd w_o(\tilde{u})$. Differentiating
$\tilde{u}_o(w_o(0,T),T)=0$ gives $-(w_o(0)/\eps)\partial_T
w_o(0)+(\partial_T\tilde{u}_o)(w_o(0),T)=0$, so $X_\sigma'(T)=w_o(0)\partial_T w_o(0)=\eps
e^{J_{o,\sigma}(T)}$. The analogous integration on $[w_i(0),\phi_\sigma(T)]$ gives $Y_\sigma'(T)=\eps
e^{-J_{i,\sigma}(T)}$.

Finally, fix $\tilde{u}_0\in[0,T]$ and put $M_0=m_{\sigma,T}(\tilde{u}_0)$. Since $\dd
w_o/\dd\tilde{u}<0$ and $w_o(s)>\phi_\sigma(s)$ on $(0,T)$ with $w_o(T)=\phi_\sigma(T)\le M_0$, we
have $w_o(\tilde{u}_0)\ge M_0$. If $w_o(\tilde{u}_0)>M_0$, let $s_*\in(\tilde{u}_0,T]$ be the first
point where $w_o(s_*)=M_0$. On $[\tilde{u}_0,s_*)$, $w_o(s)>M_0\ge\phi_\sigma(s)$, whence
\[
  \frac{\dd}{\dd s}\bigl(w_o(s)-M_0\bigr)^2 = -2\eps\,\frac{w_o(s)-M_0}{w_o(s)-\phi_\sigma(s)} \ge -2\eps.
\]
Integrating across $[\tilde{u}_0,s_*]$ yields $(w_o(\tilde{u}_0)-M_0)^2\le 2\eps(s_*-\tilde{u}_0)\le
2\eps(T-\tilde{u}_0)$. Applying the same argument to $-w_i$ and $-\phi_\sigma$ proves the second
estimate in \eqref{eq:envelope-bounds}.
\end{proof}


\subsection{Coercive tail limits, tail balancing, and inward graphs from infinity}\label{subsec:tail-inward}

For even degree $d>2$ and positive leading coefficient, both half-profiles
grow at least as fast as a positive multiple of $r^d$. A returning
half-arc then determines a finite limiting lower intercept $\ell_\sigma$
as the launch height grows. When the primitive starts in degree at least
two, one can also prescribe a finite ordinate at infinity and integrate
inward. The resulting intercept $G_\sigma(h)$ varies strictly with this
ordinate. Both assertions use the integrability of $r^{1-d}$ at infinity.


\begin{lemma}[Tail limits, strict odd monotonicity, and inward graphs]\label{lem:tail-inward}
Let $P$ be a real polynomial with $P(0)=0$ such that $P_\sigma(r)=P(\sigma r)$ satisfies
$P_\sigma(r)\ge c_0 r^d$ for all $r\ge K_0$, where $d>2$ and $c_0,K_0>0$. Suppose that either
$b_\sigma(\beta_0)<0$ exists transversely at some $\beta_0>0$, or $\min\supp(P)\ge 2$.
\begin{enumerate}
  \item[\textup{(i)}] Transverse returns $b_\sigma(\beta)<0$ exist for all $\beta\ge\beta_0$ (and for
  all $\beta>0$ when $\min\supp(P)\ge 2$), and the limit
  $\ell_\sigma(P)=\lim_{\beta\to+\infty}b_\sigma(\beta)$ is finite and strictly negative.
  \item[\textup{(ii)}] If polynomials $P_n$ with $P_n(0)=0$ converge to $P$ in $C^1$ on compact
  subsets of $\mathbb{R}$ and satisfy the uniform tail barrier $P_n(\sigma r)\ge c_0 r^d$ for all
  $r\ge K_0$, then $b_{P_n(\sigma\cdot)}(\beta_n)\to\ell_\sigma(P)$ whenever $\beta_n\to+\infty$; in
  particular, $\ell_\sigma(P_n)\to\ell_\sigma(P)$.
  The same conclusion for the finite intercepts holds under a truncated
  bound. Assume $P_n(\sigma r)\ge c_0r^d$ only for $K_0\le r\le R_n$,
  with $R_n\to\infty$. Let an actual half-arc have upper hit $\beta_n\to\infty$,
  turn $t_n\le R_n$ and transverse lower hit $b_n<0$. Then $b_n\to\ell_\sigma(P)$.
  Here the terminal ordinate of the inward leg is $P_n(\sigma t_n)$.
  The rotation argument in (i) supplies both hits from a prescribed
  finite turn when $\min\supp(P_n)\ge2$.
  \item[\textup{(iii)}] For every integer $1\le m<d$, the map $\lambda\mapsto\ell_\sigma(P+\lambda
  x^m)$ is continuous and strictly increasing (when $\sigma=+$) or strictly decreasing (when
  $\sigma=-$ and $m$ is odd) on any interval of $\lambda$ where transverse returns exist.
  Consequently, for odd $1\le m<d$, the tail gap $\Delta(P+\lambda x^m)=\ell_+(P+\lambda
  x^m)-\ell_-(P+\lambda x^m)$ increases strictly with $\lambda$.
  \item[\textup{(iv)}] If $\min\supp(P)\ge 2$, then for every $h\in\mathbb{R}$ there exists a unique
  solution $y_{\sigma,h}\in C^1([0,\infty))$ of the inward graph equation
  \begin{equation}\label{eq:inward-ode}
    y'(r) = \frac{r}{P_\sigma(r) - y(r)}, \qquad y(r) < P_\sigma(r) \quad (r\ge 0), \qquad \lim_{r\to+\infty} y(r) = h.
  \end{equation}
  Its axis intercept $G_\sigma(h)=y_{\sigma,h}(0)$ is smooth and strictly increasing on $\mathbb{R}$,
  satisfies $G_\sigma(h)<\ell_\sigma(P)$ for all $h\in\mathbb{R}$,
  $\lim_{h\to+\infty}G_\sigma(h)=\ell_\sigma(P)$, and $\lim_{h\to-\infty}(G_\sigma(h)-h)=0$. Moreover,
  under the global barrier hypotheses of \textup{(ii)} with $\min\supp(P_n)\ge 2$, one has
  $G_{\sigma,P_n}(h)\to G_{\sigma,P}(h)$ uniformly on compact $h$-intervals.
\end{enumerate}
\end{lemma}


\begin{proof}
(i) When $\min\supp(P)\ge 2$, the linearization of \eqref{eq:half-system} at $(0,0)$ has polar angular
speed $-1+O(|(r,y)|)$, bounded away from zero near $(0,0)$, so no half-orbit can approach $(0,0)$ from
$\{r>0\}$. Starting from any turn $(t,P_\sigma(t))$ ($t>0$), $r$ decreases away from $t$ in both time
directions, cannot stall at a positive radius (since $\dot{y}=-r$), stays bounded in $y$ by the envelope estimate
\eqref{eq:envelope-bounds} on each partial leg, and avoids $(0,0)$, hence hits both $\{r=0,y>0\}$ and $\{r=0,y<0\}$
transversely. This argument from a finite turn needs no assumption at
infinity and also applies to the terminal folds below. When $b_\sigma(\beta_0)<0$ is assumed, for every $\beta>\beta_0$ coercivity
$P_\sigma(r)>\beta$ forces a finite turn, and the trajectory from $(0,\beta_0)$ separates $(0,0)$ from
the inward leg from $(0,\beta)$, ensuring a transverse return $b_\sigma(\beta)<b_\sigma(\beta_0)<0$.

On the inward leg $Y'(r)=r/(P_\sigma(r)-Y(r))$ ($Y<P_\sigma$), choose $A_1\ge K_0$ so large that the
barrier $B(r)=\frac{1}{2}c_0 r^d$ satisfies $B'(r)=\frac{d}{2}c_0 r^{d-1} > \frac{2}{c_0}r^{1-d}$ for
all $r\ge A_1$ (using $d>2$). Choose $\alpha_*<0$ so negative that the solution $Y_*(r)$ of
$Y'=r/(P_\sigma(r)-Y)$ with $Y_*(0)=\alpha_*$ stays below $\min_{[0,A_1]}P_\sigma-1$ on $[0,A_1]$ and
satisfies $Y_*(A_1)<B(A_1)$. On $[A_1,\infty)$, at any hypothetical first contact $Y_*(r)=B(r)$ one
would have $Y_*'(r)=r/(P_\sigma(r)-\frac{1}{2}c_0 r^d)\le 2r^{1-d}/c_0 < B'(r)$, a contradiction.
Hence $Y_*(r)<B(r)<P_\sigma(r)$ for $r\ge A_1$, and $Y_*<P_\sigma$ also on $[0,A_1]$, so $Y_*$ never
turns. Scalar comparison of $Y'=r/(P_\sigma(r)-Y)$ shows $b_\sigma(\beta)>\alpha_*$ for all $\beta$,
and monotonicity gives $\ell_\sigma(P)\in[\alpha_*,b_\sigma(\beta_0)]$.

(ii) Fix $\alpha<\ell_\sigma(P)$ and put $\delta=\ell_\sigma(P)-\alpha>0$. As $\beta\to+\infty$, the
turns $t(\beta)$ of returning orbits diverge by \eqref{eq:envelope-bounds}, and their inward legs
$Y_\beta^-(r)$ satisfy $Y_\beta^-(0)=b_\sigma(\beta)>\ell_\sigma(P)=\alpha+\delta$. Since vertical
gaps between solutions of $Y'=r/(P_\sigma(r)-Y)$ are nondecreasing in $r$ (as
$\partial_Y[r/(P_\sigma-Y)]>0$), the solution $Y_\alpha$ from $Y_\alpha(0)=\alpha$ exists for all
$r\ge 0$ and satisfies $P_\sigma(r)-Y_\alpha(r)>Y_\beta^-(r)-Y_\alpha(r)\ge \delta$, whence
$Y_\alpha'(r)\le r/\delta$ and $Y_\alpha(r)\le \alpha+r^2/(2\delta)$. Choose $A_1\ge K_0$ as above
with $\alpha+A_1^2/(2\delta)<\frac{1}{2}c_0 A_1^d$. By compact $C^1$ convergence on $[0,A_1]$, for all
large $n$ the solution $Y_{\alpha,n}$ from $\alpha$ remains strictly below
$P_n(\sigma r)$ on $[0,A_1]$ and satisfies $Y_{\alpha,n}(A_1)<B(A_1)$. The first-contact argument in (i) extends it
strictly below $B$ up to $R_n$ in the truncated case. The turns
$t_n$ tend to infinity: on a bounded radial interval the upper-envelope
bound and compact convergence bound the upper hit. Thus
\[
 Y_{\alpha,n}(t_n)<B(t_n)<P_n(\sigma t_n)=Y_n^-(t_n).
\]
Ordering just before the turn and uniqueness back to the axis give
$b_n>\alpha$. For the reverse bound fix a returning launch $B_0>\beta_0$
of the limit profile. Its transverse arc persists on a compact set for
$P_n$, with $b_{P_n(\sigma\cdot)}(B_0)\to b_\sigma(B_0)$.
The arc from $\beta_n>B_0$ cannot cross this inner arc, so its lower
hit lies below that from $B_0$. Taking the lower and upper limits, then
$\alpha\uparrow\ell_\sigma(P)$ and $B_0\to\infty$, proves the claim.
For the global barrier all sufficiently high launches turn and return,
by (i) and persistence of the fixed inner arc. Choosing those launches
arbitrarily far out also gives convergence of the limiting intercepts.
In the truncated assertion the comparison at $t_n$ uses the turning
ordinate itself. Prescribing a bounded ordinate $h$ at a distant cut
instead leads to the incoming graph of (iv).

(iii) Since $1\le m<d$, $(P+\lambda x^m)(\sigma r)\ge \frac{1}{2}c_0 r^d$ holds on $[K_1,\infty)$
locally uniformly in $\lambda$, so continuity of $\lambda\mapsto\ell_\sigma(P+\lambda x^m)$ follows
from~(ii). Take $\lambda_2>\lambda_1$ with $\sigma=+$, so $Q_2(r)=P(r)+\lambda_2 r^m > P(r)+\lambda_1
r^m = Q_1(r)$ for all $r>0$. For every sufficiently large common
launch both returns exist by (i). Lemma~\ref{lem:variational}(ii) gives $b_{Q_2}(\beta)>b_{Q_1}(\beta)$,
hence $\ell_+(Q_2)\ge \ell_+(Q_1)$. Chen, Jin, Li and Zhou prove the
strict order for their trinomial profiles by a uniform parameter
variation \cite[Lemma~2.9, (2.41)--(2.45)]{CJLZ}. Here it follows
by comparing the two limiting graphs on a fixed short interval. Suppose toward a contradiction that
$\ell_+(Q_2)=\ell_+(Q_1)=\ell<0$. Choose $r_0>0$ so small that $\min_{[0,r_0]} Q_i \ge \ell/2 > \ell$
for $i=1,2$. On $[0,r_0]$, the ODEs $Y'=r/(Q_i(r)-Y)$ are smooth and uniformly Lipschitz on $\{Y\le
3\ell/4\}$. Since $b_{Q_i}(\beta)\downarrow\ell$, continuous dependence on $[0,r_0]$ implies
$Y_{i,\beta}^-(r)\to Y_i(r)$ uniformly on $[0,r_0]$ as $\beta\to+\infty$, where $Y_i$ solves
$Y_i'=r/(Q_i-Y_i)$ with $Y_i(0)=\ell$. Because $Q_2(r)>Q_1(r)$ on $(0,r_0]$ and $Y_1(0)=Y_2(0)=\ell$,
differential comparison gives $Y_2(r_0)<Y_1(r_0)$. Yet at every large finite $\beta$, integrating
\eqref{eq:filippov} downward from $u_{1,\beta}(\beta)=u_{2,\beta}(\beta)=0$ shows
$u_{2,\beta}(y)<u_{1,\beta}(y)$ for all $y\in[b_{Q_2}(\beta),\beta)$; because the inward legs
$Y_{i,\beta}^-(r)$ increase strictly with $r$, lying strictly inside the $Q_1$-arc means
$Y_{2,\beta}^-(r_0)>Y_{1,\beta}^-(r_0)$. Passing to $\beta\to+\infty$ yields $Y_2(r_0)\ge Y_1(r_0)$, a
contradiction.

(iv) Fix $h\in\mathbb{R}$ and choose $K_h\ge K_0$ such that $P_\sigma(r)\ge c_0 r^d >
\max(h,0)+2|h|+2$ for all $r\ge K_h$. For each $K\ge K_h$, let $y_K(r)$ solve
$y'(r)=r/(P_\sigma(r)-y(r))$ backwards on $[0,K]$ from $y_K(K)=h$. Because $\min\supp(P)\ge 2$, the
ordinary returning arc turning at $(K,P_\sigma(K))$ has an inward leg $Y_{(K)}^-(r)$ defined on
$[0,K]$ with $Y_{(K)}^-(r)<P_\sigma(r)$ on $[0,K)$ and $Y_{(K)}^-(K)=P_\sigma(K)>h=y_K(K)$. By
uniqueness for the scalar ODE on $[0,K]$, $y_K(r)<Y_{(K)}^-(r)<P_\sigma(r)$ throughout $[0,K]$ and
$y_K'(r)>0$, so $y_K(r)\le h$ on $[0,K]$. For $K_2>K_1\ge K_h$,
\[
  0 \le h - y_{K_2}(K_1) = \int_{K_1}^{K_2} \frac{r\dd r}{P_\sigma(r)-y_{K_2}(r)} \le
  \int_{K_1}^\infty \frac{r\dd r}{P_\sigma(r)-h} \xrightarrow{K_1\to\infty} 0.
\]
For any two solutions $y_a,y_b<P_\sigma$ on $[0,K_1]$, $\frac{\dd}{\dd
r}|y_a(r)-y_b(r)|=\frac{r|y_a-y_b|}{(P_\sigma-y_a)(P_\sigma-y_b)}>0$, so backward integration from
$K_1$ to $0$ strictly contracts $|y_a-y_b|$. Hence $\{y_K\}$ is uniformly Cauchy on $[0,K_1]$ for
every $K_1\ge K_h$, converging to a solution $y_{\sigma,h}$ of \eqref{eq:inward-ode}. Uniqueness
follows because any two solutions with limit $h$ have $|y_1(r)-y_2(r)|$ nondecreasing on $[0,\infty)$
and tending to $0$ at $+\infty$.

Comparing $y_{\sigma,h}$ on $[0,t(\beta)]$ with $Y_\beta^-$ for $P_\sigma(t(\beta))>h$ gives
$G_\sigma(h)<b_\sigma(\beta)$, so $G_\sigma(h)\le \ell_\sigma(P)$; strict monotonicity of $h\mapsto
G_\sigma(h)$ (two graphs of distinct limiting heights are strictly ordered at a sufficiently large
finite cut, and uniqueness preserves strict order to the axis) then yields
$G_\sigma(h)<G_\sigma(h+1)\le \ell_\sigma(P)$. Conversely, in part~(ii) we showed that for every
$\alpha<\ell_\sigma(P)$, the inward solution $Y_\alpha$ from $Y_\alpha(0)=\alpha$ satisfies
$Y_\alpha(r)\le\alpha+r^2/(2\delta)<\frac{1}{2}c_0 r^d$ for large $r$, whence
$Y_\alpha'(r)=O(r^{1-d})$ integrates to a finite limit $h(\alpha)=\lim_{r\to\infty}Y_\alpha(r)$, so
$G_\sigma(h(\alpha))=\alpha$ and $\lim_{h\to+\infty}G_\sigma(h)=\ell_\sigma(P)$. For $h<M_*-1$ where
$M_*=\min_{r\ge 0}P_\sigma(r)$, $0\le h-G_\sigma(h)=\int_0^\infty\frac{r\dd
r}{P_\sigma(r)-y_{\sigma,h}(r)}\le \int_0^\infty\frac{r\dd r}{P_\sigma(r)-h}\to 0$ as $h\to-\infty$ by
dominated convergence. Differentiation at a finite cutoff followed by passage to the limit gives
$G_\sigma'(h)=\exp[-\int_0^\infty r(P_\sigma-y_{\sigma,h})^{-2}\dd r]>0$.
The differentiated tail converges uniformly for $h$ in a compact interval
by the same polynomial lower bound; higher derivatives follow by successive
differentiation of the graph equation. Finally, under uniform coercivity $P_n(\sigma r)\ge c_0 r^d$
($r\ge K_0$), the tail loss $\int_{K_1}^\infty r/(c_0 r^d-h)\dd r$ is uniform on compact $h$-intervals
and backward contraction transfers compact convergence on $[0,K_1]$ to $G_{\sigma,P_n}(h)\to
G_{\sigma,P}(h)$.
\end{proof}


\subsection{Simultaneous hyperbolicity on a fixed support}\label{subsec:hyperbolicity}

De Maesschalck and Dumortier observed that each odd coefficient of
the classical primitive is a rotation parameter \cite[pp.~2163--2164]{DMD}.
Indeed, for $X_\mu=(y-F(x)-\mu x^m,-x)$ with $m$ odd,
$\det(X_\mu,\partial_\mu X_\mu)=-x^{m+1}$ has one sign off the axis.
Cai, Chen, Jin and Lu use a linear rotation and a regular value to
obtain simple fixed points simultaneously for a continuous piecewise
linear primitive \cite[proof of Theorem~1.1 in \S2]{CCJL}. Lu uses
the linear coefficient for polynomial returns
\cite[Proposition~2.4(i)]{LuDegree}. In our fixed support we can use
any odd exponent already present. Formula~\eqref{eq:b-param} gives
strict parameter derivatives of opposite signs for the two half-returns.
This allows all selected brackets to be regularized with a single change
of coefficient.

\begin{lemma}[Single-field regularization and openness]
Let\label{lem:hyperbolicity} $F_0(x)=\sum_{i=1}^k a_i x^{m_i}$ ($a_i\ne 0$) contain an odd exponent $m_*=m_{i_0}\in\supp(F_0)$,
and set $F_\mu(x)=F_0(x)+\mu x^{m_*}$. Suppose that at $\mu=0$:
\begin{enumerate}
  \item[\textup{(a)}] $I_1,\dots,I_N\subset(0,\infty)$ are pairwise disjoint compact intervals
  throughout which both transverse half-returns $b_{\pm,0}(\beta)<0$ exist, and
  $\oscD_0(\beta)=b_{+,0}(\beta)-b_{-,0}(\beta)$ has opposite strict signs at the two endpoints of
  each $I_r$ ($1\le r\le N$);
  \item[\textup{(b)}] $\beta_1^*,\dots,\beta_M^*\in(0,\infty)\setminus\bigcup_{r=1}^N I_r$ are
  distinct hyperbolic periodic-orbit intercepts ($\oscD_0(\beta_s^*)=0$, $\oscD_0'(\beta_s^*)\ne 0$).
\end{enumerate}
Then there exists $\mu_0\in(0,|a_{i_0}|)$ such that for Lebesgue-almost-every $\mu\in(-\mu_0,\mu_0)$,
the polynomial $F_\mu$ satisfies $\supp(F_\mu)=\supp(F_0)$ and system \eqref{eq:lienard} for $F=F_\mu$
admits at least $N+M$ distinct hyperbolic limit cycles. Consequently, the set of coefficient vectors
on $\supp(F_0)$ realizing at least $N+M$ distinct hyperbolic limit cycles has nonempty interior in
$(\mathbb{R}\setminus\{0\})^k$.
\end{lemma}


\begin{proof}
Choose disjoint compact intervals $J_1,\dots,J_M\subset(0,\infty)\setminus\bigcup_{r=1}^N I_r$ with
$\beta_s^*\in\operatorname{int}(J_s)$ on which $\oscD_0'\ne 0$ and $\oscD_0$ has opposite strict signs
at $\partial J_s$. By real-analytic dependence of transverse half-returns and the implicit function
theorem on $J_1,\dots,J_M$, there exists $\mu_0\in(0,|a_{i_0}|)$ such that for all
$\mu\in(-\mu_0,\mu_0)$: (i) $a_{i_0}+\mu\ne 0$, so $\supp(F_\mu)=\supp(F_0)$; (ii)
$b_{\pm,\mu}(\beta)<0$ remain well-defined and real analytic on an open neighborhood $U$ of
$\bigl(\bigcup_{r=1}^N I_r\bigr)\cup\bigl(\bigcup_{s=1}^M J_s\bigr)$; (iii)
$\oscD_\mu(\beta)=b_{+,\mu}(\beta)-b_{-,\mu}(\beta)$ retains its strict opposite endpoint signs on
every $I_r$; and (iv) each $J_s$ contains a unique zero $\beta_s^*(\mu)$ of $\oscD_\mu$ with
$\oscD_\mu'(\beta_s^*(\mu))\ne 0$.

Because $m_*$ is odd, $P_{\sigma,\mu}(r)=P_{\sigma,0}(r)+\sigma\mu r^{m_*}$ has $\partial_\mu
P_{\sigma,\mu}(r)=\sigma r^{m_*}$. By \eqref{eq:b-param} in Lemma~\ref{lem:variational}(ii), since
$-b_{\sigma,\mu}>0$ and $u_\sigma(y)>0$ on $(b_{\sigma,\mu},\beta)$, one has $\partial_\mu
b_{+,\mu}(\beta)>0$ and $\partial_\mu b_{-,\mu}(\beta)<0$, whence
\[
  \partial_\mu \oscD_\mu(\beta) = \partial_\mu b_{+,\mu}(\beta) - \partial_\mu b_{-,\mu}(\beta) > 0
  \qquad \text{for all } (\beta,\mu)\in\Omega := \Bigl(\bigcup_{r=1}^N
  \operatorname{int}(I_r)\Bigr)\times(-\mu_0,\mu_0).
\]
By the implicit function theorem, $\mathcal{Z}=\{(\beta,\mu)\in\Omega : \oscD_\mu(\beta)=0\}$ is the
graph $\mu=\mathcal{M}(\beta)$ of a real-analytic function $\mathcal{M}:W\to(-\mu_0,\mu_0)$ defined on
an open subset $W\subset\bigcup_{r=1}^N\operatorname{int}(I_r)$. Differentiating
$\oscD_{\mathcal{M}(\beta)}(\beta)=0$ gives
$\partial_\beta\oscD_{\mathcal{M}(\beta)}(\beta)+\partial_\mu\oscD_{\mathcal{M}(\beta)}(\beta)\,\mathcal{M}'(\beta)=0$.
At a zero $\beta\in\operatorname{int}(I_r)$, degeneracy
($\partial_\beta\oscD_\mu(\beta)=0$, equivalently $\Pi_\mu'(\beta)=1$
by Lemma~\ref{lem:variational}(ii)) is therefore equivalent to
$\mathcal M'(\beta)=0$. Thus a parameter admits a degenerate zero in
these intervals precisely when it is a critical value of $\mathcal M$. By Sard's
theorem, the critical values of $\mathcal{M}$ have measure zero in $(-\mu_0,\mu_0)$. Choosing any
regular value $\mu\in(-\mu_0,\mu_0)$, the intermediate value theorem supplies at least one zero of
$\oscD_\mu$ in each $\operatorname{int}(I_r)$ ($1\le r\le N$), every such zero satisfies
$\partial_\beta\oscD_\mu(\beta)\ne 0$ (hence $\Pi_\mu'(\beta)\ne 1$), and these $N$ hyperbolic limit
cycles are disjoint from one another and from $\beta_1^*(\mu),\dots,\beta_M^*(\mu)$. At this fixed
$\mu$, all $N+M$ hyperbolic fixed points of $\Pi_\mu$ persist smoothly and remain distinct and
hyperbolic on an open neighborhood of the coefficient vector of $F_\mu$ in
$(\mathbb{R}\setminus\{0\})^k$.
At a marked zero the same formula \eqref{eq:b-param} gives
$\partial_\mu\oscD_\mu>0$. Its launch moves with derivative
$(\beta_s^*)'(\mu)=-\partial_\mu\oscD_\mu/\partial_\beta\oscD_\mu\ne0$.
\end{proof}

\begin{remark}[Local cyclicity on a fixed support]\label{rem:local-cyclicity}
Fix an even primitive and allow the coefficients on a prescribed
set of $b\ge1$ odd exponents $m'_1<\cdots<m'_b$ to vary near zero.
We count only cycles shrinking to this reversible center. In the
coordinates $(x,y-F(x))$ the damping in Fran\c{c}oise's convention
is $p=-F'$. At zero trace, Theorem~1.1 and Lemmas~1.3--1.4 of the
published paper \cite[pp.~257--259]{Francoise} identify its even
damping coefficients as generators of the Bautin ideal and give
their triangular order in the inverse return series.

Write $c_i$ for the coefficient of $x^{m'_i}$ and $\rho$ for the
radius on a local polar section. Suppress the even parameters, which
may vary near their fixed values. Let $D_i$ be the inverse return
displacement with $c_1=\cdots=c_{i-1}=0$, and put $D_{b+1}=0$
by reversibility. For $m'_i\ge3$, the triangular identities make
every coefficient below order $m'_i$ vanish in $D_i$ and give a
nonzero constant times $c_i$ at that order. If $m'_1=1$, the linear
term follows from the inverse focus multiplier, whose derivative in
$c_1$ at zero is nonzero. In both cases the analytic Hadamard
factorization in $c_i$ is
\[
 D_i-D_{i+1}
 =c_i\int_0^1\partial_{c_i}D_i(\rho,t c_i,c_{i+1},\ldots,c_b)\,\dd t
 =c_i\rho^{m'_i}u_i(\rho,c).
\]
The coefficient statements give analytic $u_i$ with $u_i(0,0)\ne0$.
Summing yields $D_1=\sum_i c_i\rho^{m'_i}u_i$. For fixed parameters
the leading coefficient of each initial Wronskian of these functions
$\rho^{m'_i}u_i$ is the Vandermonde product of the corresponding
exponents times the product of their $u_i(0,c)$. All these coefficients
stay away from zero on a sufficiently small closed parameter
neighborhood. The Wronskians therefore have fixed nonzero signs for
small $\rho>0$, uniformly there. Successive applications of Rolle's
theorem give at most $b-1$ small isolated cycles. This local
fixed-support deduction concerns one center. The separated finite
orbits counted by Theorem~\ref{thm:main} occur at different radii.
\end{remark}
